\section{Conclusion}
\label{section:6.1}

This thesis has presented \textbf{FreelanceChain}, a proof-of-concept decentralized freelance marketplace implemented on Solana Devnet through the Anchor framework. The prototype addresses a clearly documented structural deficiency of the contemporary online freelance industry -- the concentration of trust, payment custody, dispute resolution, identity verification, and governance within a single private intermediary -- by replacing several of those centralized functions with smart-contract logic. The work demonstrates technical feasibility on Solana Devnet at the prototype level; the broader claim that a decentralized marketplace can replace incumbent platforms in production remains an open empirical question to be settled by future Mainnet deployment, audit, and user study.

Relative to \textbf{Upwork and Fiverr}, FreelanceChain demonstrates how on-chain escrow, on-chain arbitration, and portable SBT reputation can be implemented without a custodial intermediary; the acknowledged trade-off is greater onboarding friction and dependence on user-managed wallets. Relative to \textbf{Ethlance}, FreelanceChain explores a broader feature surface within a single Anchor program -- milestone escrow, biometric KYC, SBT reputation, DAO governance, ZK credential storage, AI risk advisory, social recovery, and an AI-oracle interface -- and uses Solana's cost profile (below USD\,0.001 per transaction at Devnet fee-market rates) to render small-value engagements plausible. Relative to \textbf{Colony and Braintrust}, FreelanceChain targets individual freelance project matching rather than team-oriented task or treasury management.

The objectives formulated in Section~\ref{section:1.2} have been achieved at the prototype level: eighteen Anchor modules covering the engagement lifecycle have been designed, implemented, and deployed to Devnet; a Next.js front-end of 53 React components has been constructed; a browser-side biometric KYC pipeline has been integrated; and 38 functional test cases have been executed without failure. Five items lie outside the present scope and are explicitly identified as limitations rather than completed work: \textbf{(i)} Mainnet deployment was not pursued; \textbf{(ii)} a third-party security audit has not been performed; \textbf{(iii)} a formal large-scale user study has not been conducted -- the survey of Section~\ref{subsection:2.1.2} is exploratory only; \textbf{(iv)} the ZK module implements commitment storage but not a full on-chain ZKP verification circuit; and \textbf{(v)} the governance module's staking-reward distribution sub-module is implemented only as a stub. Active liveness detection in the KYC pipeline and a larger empirical benchmark for the AI advisory subsystem are likewise deferred.

Several methodological observations may be extracted. Solana's strict separation of code and state requires greater up-front design discipline than Ethereum's combined model but yields programs that are more parallelizable and cheaper to query. Anchor's constraint macros help eliminate a class of common vulnerabilities, though careful attention is required to the ordering of validation checks in complex multi-account instructions. Browser-side machine learning is feasible for initial-implementation biometric verification provided that model size is managed carefully and that the absence of active liveness checks is acknowledged as a limitation. Integration of an LLM API into a Web3 application requires deliberate engineering of the trust boundary: keeping AI output strictly advisory -- never committed on-chain -- prevents the AI subsystem from becoming a new centralized point of failure.

\section{Future Work}
\label{section:6.2}

Six directions are identified as priorities for future development.

\textbf{Mainnet deployment and security audit.} Solana Mainnet release requires careful economic modeling of the governance-token supply, DAO-treasury funding, and referral-commission distribution, together with a regulatory-compliance analysis concerning the legal classification of the governance token. A third-party security audit of the Anchor program is a prerequisite to any custody of real funds.

\textbf{Full on-chain ZKP verification.} The present ZK module stores commitments without verifying Groth16 or PLONK proofs on-chain. Client-side proof generation through Bellman or Gnark, paired with on-chain verification (potentially through the Light Protocol stack), would enable fully trustless credential proofs and eliminate the dependence on oracle nodes.

\textbf{Decentralized dispute arbitration.} A staking-based jury-selection mechanism inspired by Kleros -- in which token holders stake to serve as arbitrators and forfeit a portion of stake when voting against the majority -- would render arbitration credibly neutral and more resistant to bribery than the present small-arbitrator-set design.

\textbf{Cross-chain identity portability.} Extension of the KYC record to support the W3C Decentralized Identifier (DID) standard would permit verified users to present their attestation to compliant Ethereum-based contracts via bridges or cross-chain messaging, without re-verification.

\textbf{Active liveness detection and a larger KYC empirical study.} Adding blink-based or head-turn liveness challenges (e.g., via MediaPipe FaceMesh) would close the still-image attack surface acknowledged in Section~\ref{section:5.3}. A demographically balanced empirical study of the threshold's true positive / false positive rate is also required to upgrade the illustrative threshold check of Table~\ref{table:kyc_threshold} into a true evaluation.

\textbf{DAO staking-reward distribution and AI subsystem hardening.} Completion of the staking-reward sub-module identified in Section~\ref{section:5.5}, a multilingual evaluation of the AI risk advisor, and stronger prompt-injection defences (structured-output enforcement, evidence-passage citation, signed-system-prompt patterns) would harden the two prototype areas left explicitly incomplete in the present work.

\textbf{Observability extensions.} The read-only Prometheus / Grafana sidecar described in Section~\ref{section:4.7} establishes a baseline view of program health but is deliberately narrow. Three extensions are tractable next steps: pushing precise submit-to-confirm latency from the dApp to a Prometheus Pushgateway (which would turn Table~\ref{table:perf} into a continuous time series rather than a one-shot harness measurement); shipping an Alertmanager rule set covering heartbeat age, failed-instruction rate, and unexpected escrow drains; and exposing the same metric set as a public read-only Grafana instance so that auditors and users can verify the operator's claims independently.

In conclusion, the FreelanceChain prototype demonstrates the technical feasibility of building a comprehensive, trust-decentralizing, and privacy-respecting online freelance marketplace on contemporary public blockchain infrastructure. The five core contributions -- dual-track escrow, SBT reputation, browser-side biometric KYC, Groq/LLaMA-powered AI advisory, and DAO governance -- collectively address several of the structural deficiencies of incumbent centralized platforms and open multiple avenues for future research at the intersection of blockchain systems, privacy-preserving computation, decentralized governance, and AI-augmented economic coordination. The work makes no claim to having solved the problem; rather, it offers a working starting point that subsequent research can refine, audit, and scale.
